\def\mymode{0} % Mode 0 for presentation, mode 1 (or not 0) for a handout.
\if 0\mymode
    \documentclass[hyperref={colorlinks = true, urlcolor = magenta}]{beamer}
\else
    \documentclass[handout]{beamer}
    \usepackage{handoutWithNotes}
    \pgfpagesuselayout{3 on 1 with notes}[a4paper]
\fi
% --- Begin Preamble --- % (fold)
\setbeamerfont{footnote}{size=\tiny}
\graphicspath{{./Images/},{../Images/}} % the folder in which images are stored for this project
\usepackage{layout,verbatim,tikz,graphicx,amsmath,amssymb,appendixnumberbeamer,adjustbox,subcaption}
% \usepackage[style=verbose]{biblatex}
\usetikzlibrary{calc,shapes,arrows}
\definecolor{DBlack}{RGB}{35,31,32}
\definecolor{DPurple}{RGB}{126,49,123}
\definecolor{DBlue}{RGB}{0,99,136}

\newtheorem*{defn}{Definition}
\theoremstyle{definition}

\newtheorem*{thm}{Theorem}
\theoremstyle{plain}

\usecolortheme[RGB={83,13,88}]{structure}
% This is the theme that we are modifying
\usetheme{Boadilla}

% Change the font colour to white on the title page to stand out against the purple
\setbeamercolor{title}{fg=white}
\setbeamerfont{title}{family=\rmfamily}
\setbeamercolor{author}{fg=white}
\setbeamerfont{author}{series=\bfseries}
\setbeamercolor{institute}{fg=white}
\setbeamercolor{date}{fg=white}
\setbeamerfont{date}{series=\bfseries}

% Remove the default ball bullet points in the toc, but keep the numbering
\setbeamertemplate{sections/subsections in toc}[sections numbered]
% Change the item icons to be circles rather than balls
\setbeamertemplate{itemize items}[circle]
% Define blocks to be rounded rectangles with shadows
\setbeamertemplate{blocks}[rounded][shadow=true]

% Turn the @ symbol into character class other
\makeatother
% Modify the footline to get a two part footline with subsection names
\setbeamertemplate{footline}
{
  \leavevmode%
  \hbox{%
  % Create the first box width 0.4 of slide, containing the Author's Name (Short)
  \begin{beamercolorbox}[wd=.4\paperwidth,ht=2.25ex,dp=1ex,center]{author in head/foot}%
    \usebeamerfont{author in head/foot}\insertshortauthor%\ (\insertinstitute)
  \end{beamercolorbox}%
  % Create the second box width 0.6 of slide, containing Subsection name, current slide and total slides
  \begin{beamercolorbox}[wd=.6\paperwidth,ht=2.25ex,dp=1ex,center]{subsection in head/foot}%
    \usebeamerfont{subsection in head/foot}\insertsubsection\hspace*{3em}
    \insertframenumber{} / \inserttotalframenumber\hspace*{1ex}
  \end{beamercolorbox}}%
  \vskip0pt%
}
% Turn the @ symbol back to a letter
\makeatletter

% Remove the default navigation symbols from the bottom of the slides
\setbeamertemplate{navigation symbols}{}

% --- Start Custom Commands --- % (fold)
\newcommand{\cfbox}[1]{\adjustbox{cfbox=DPurple 1pt}{#1}}
% --- End Custom Commands --- % (end)

% --- End Preamble --- % (end)

\author{Sam Fearn}
\institute[Durham]{Durham University}
\title{A practical beginner guide to the STACK API}
\date{October 22\textsuperscript{nd}, 2026}

% The following code is inserted at every new subsection. It adds an outline slide with the current section and subsection highlighted.
% \AtBeginSection[]
% {
%    \begin{frame}[plain,noframenumbering]
%        \frametitle{Outline}
%        \tableofcontents[currentsection,currentsubsection]
%    \end{frame}
% }

\begin{document}

% --- Title and Contents --- % (fold)
% --- The Title Slide --- %
{
% Draw a purple box for the background of the title slide and add the durham logo at the bottom
\setbeamertemplate{background canvas}{
\begin{tikzpicture}
    \clip (0,0) rectangle (\paperwidth,\paperheight);
    \fill[color=DPurple] (0,0) rectangle (\paperwidth,\paperheight);
    \node[inner sep=0pt] (MathsLogo) at (1.5,8.3)
        {\includegraphics[width=20mm]{MathsLogo.png}};
    \node[inner sep=0pt] (DurhamLogo) at (2.5,1)
        {\includegraphics[width=50mm]{DULogo_WHITE.png}};
\end{tikzpicture}
}

\begin{frame}[plain,noframenumbering]
\maketitle
\end{frame}
}

%--- Complete Outline --- %
% \begin{frame}[plain,noframenumbering]
%        \frametitle{Outline}
%        \tableofcontents
% \end{frame}
%(end)

\section{Overview} % (fold)
\label{sec:overview}
\begin{frame}{Aims}
    \begin{itemize}
        \item Explain the basics of what the STACK API is/does.
        \item Workshop converting a sample tex file into html with an embedded example stack question (follow-along).
        \item We will use \href{https://www.docker.com/products/docker-desktop/}{Docker Desktop} (\url{https://www.docker.com/products/docker-desktop/}).
    \end{itemize}
    \medskip
    \begin{center}\href{https://www.docker.com/products/docker-desktop/}{\includegraphics[width=0.6\textwidth]{docker-logo}}\end{center}
\end{frame}

% section overview (end)

\section{Formative Assessment} % (fold)
\label{sec:formative_assessment}

\begin{frame}[t]{Moodle Quizzes}
    \url{https://www.docker.com/products/docker-desktop/}
    \begin{center}
        \includegraphics[width=0.5\textwidth]{Moodle-logo}
    \end{center}
    
    STACK is an excellent tool for automating mathematics assessment, and as a Moodle quiz question plugin instructors can:
    \begin{itemize}
        \item<2-> Use STACK for formal assessments, with exportable grade data,
        \item<3-> Use question behaviours to control how students interact with the questions,
        \item<4-> Control timed release windows, with user/group overrides,
        \item<5-> Use activity completion settings to further control access.
    \end{itemize}
    
    \pause\medskip
    
    \uncover<6->{Most of these features are likely necessary for \emph{summative} assessment (assessment \emph{of} learning), but may not be for \emph{formative} assessment (assessment \emph{for} learning).}
\end{frame}

\begin{frame}[t]{Formative engagement barrier?}
    \url{https://www.docker.com/products/docker-desktop/}
    \only<1>{\begin{center}\includegraphics[width=0.75\textwidth]{Ultra1}\end{center}}
    \only<2>{\begin{center}\includegraphics[width=0.75\textwidth]{Ultra2}\end{center}}
    \only<3>{\begin{center}\includegraphics[width=0.75\textwidth]{Moodle1}\end{center}}
    \only<4>{\begin{center}\includegraphics[width=0.75\textwidth]{Moodle2}\end{center}}
    \only<5>{\begin{center}\includegraphics[width=0.75\textwidth]{Ultra3}\end{center}}
    \only<6>{\begin{center}\includegraphics[width=0.75\textwidth]{Ultra4}\end{center}}
    \only<7>{\begin{center}\includegraphics[width=0.75\textwidth]{ODENotes}\end{center}}
\end{frame}

% section formative_assessment (end)

\section{The STACK API} % (fold)
\label{sec:the_stack_api}
\begin{frame}{Questions in notes}
    \begin{center}
        \includegraphics[width=0.65\textwidth]{no-moodle}
    \end{center}
    The STACK API allows questions to be used outside of Moodle, allowing for the possibility of embedding STACK questions directly inside html lecture notes.
    
    \medskip
    
    \begin{center}
        \huge{DEMO}
    \end{center}
    
\end{frame}

\begin{frame}[t]{Setup}
    \begin{center}
        \includegraphics[width=0.2\textwidth]{docker-logo}
    \end{center}
    The STACK API can be installed and run via Docker.
    \only<1>{\begin{center}\includegraphics[width=0.8\textwidth]{dockerdesktop}\end{center}}
    
    We can then communicate with the API at localhost:3080.
    \medskip
    
    \uncover<2->{The API provides 5 routes:}
    \begin{itemize}
        \item<2-> (localhost:3080)/render
        \item<3-> /validate
        \item<4-> /grade
        \item<5-> /download
        \item<6-> /test
    \end{itemize}
    
    \medskip
    \uncover<7->{Each route expects to receive a POST request, whose body is a JSON object containing the question XML (and possibly other fields).}
    
    \medskip
    \uncover<7->{Each route responds with a JSON object.}
\end{frame}
% section the_stack_api (end)

\section{Creating dynamic notes} % (fold)
\label{sec:creating_dynamic_notes}

\begin{frame}[t,fragile]{Embedding questions in notes}
    \begin{itemize}
        \item<1-> The notes shown in the demo were originally written in \LaTeX{}.
        \item<2-> A new markdown source file was created using Pandoc, plus some custom extensions. \uncover<3->{The mathematics (inline and display) remains in \LaTeX{}.}
        \item<4-> This markdown can then be processed (again, Pandoc+) into both html and pdf outputs \uncover<5->{(can also go straight to html from \LaTeX{} if preferred).}
        \item<6-> The required questions should be exported from Moodle in Moodle XML format.
        \item<7-> Question blocks are included in the markdown using the following syntax:\\
    \hspace{2em}\verb|::: {.que .stack data-qid="XXXX"}:::| \\
    where \verb|XXXX| is the Question Name of the required question contained in the Moodle XML.
        
        \only<8>{If going straight to html, the following can instead be inserted into the html directly:\\
    \hspace{2em}\texttt{<div class="que stack" data-qid="XXXX''></div>}.}
        \item<9->The html notes and question XML file/s (as well as an additional `stackapicalls.js') are hosted on an appropriate webserver.
    \end{itemize}    
\end{frame}

\begin{frame}[t]{Process - key steps}
    \only<1-6>{The notes load some JS, based on the code used in the example `(localhost:3080)/stack.php' page.
    
    \begin{itemize}
        \item<2-> The (empty) divs of class ``que stack'' from before are numbered and populated, initially containing a button to generate a question.
        \only<2>{\begin{center}\includegraphics[width=0.9\textwidth]{API1}\end{center}}
        \item<3-> This button gets the XML of the question from a preconfigured list of XML files, and passes this to the /render route, before rendering the response.
        \only<3>{\begin{center}\includegraphics[width=0.9\textwidth]{API2}\end{center}}
        \item<4-> Event listeners are attached to the inputs to call the /validate route when the input changes.
        \only<4>{\begin{center}\includegraphics[width=0.9\textwidth]{API3}\end{center}}
        \item<5-> A random seed is used for the render, so the seed is attached to the answer submit button.
        \item<6-> The answer submit button calls the /grade route, passing the xml as before, with the matching seed to the /render call. The response from this is then rendered to the page.
    \end{itemize}}
    \only<7>{\begin{center}\includegraphics[width=0.9\textwidth]{API4}\end{center}}
\end{frame}

\begin{frame}[t]{Notes and caveats}
    \begin{itemize}
        \item<1-> Currently, the routes expect the question xml to be sent as part of the request. This means a knowledgeable user can obtain the full question xml. \uncover<2->{This may not be acceptable if the same questions appear in formal assessments.}
        \item<3-> Similarly, the response from /render contains the model answer. \uncover<4->{This is useful for a `reveal answer/s', even if the general/specific feedback doesn't contain this.}
        \item<5-> In the example, the questions are essentially always in `Adaptive mode'.
    \end{itemize}
\end{frame}

% section creating_dynamic_notes (end)

\section{Conclusions} % (fold)
\label{sec:conclusions}
\begin{frame}[t]{Conclusions}
    \begin{itemize}
        \item<1-> html lecture notes are more accessible for all users, and should be standard.
        \item<2-> In addition, they provide possibilities for enhanced interactivity, for example by embedding STACK questions directly into the notes.
        \item<3-> This in turn enables closer integration of course materials, which I expect will lead to students making better use of practice questions.
        \item<4-> We can also embed other interactive elements: \only<4->{Desmos}\only<5->{, Geogebra}\only<6->{, JupyterLite}\only<7->{, Visual PDE}\only<8->{, \ldots}
        \only<4>{\begin{center}\includegraphics[width=0.5\textwidth]{Desmos}\end{center}}
        \only<5>{\begin{center}\includegraphics[width=0.55\textwidth]{Geogebra}\end{center}}
        \only<6>{\begin{center}\includegraphics[width=0.6\textwidth]{JupyterLite}\end{center}}
        \only<7>{\begin{center}\includegraphics[width=0.6\textwidth]{VisualPDE}\end{center}}
        \item<9-> While formal assessment questions are good to include, there is also (massive) scope for additional questions (and proofs), and more scaffolded examples.
    \end{itemize}
\end{frame}

% section conclusions (end)

\appendix
%Slides in this section won't count towards the numbering of the main talk. They will instead have their own count.

\section*{Thanks} % (fold)
{\setbeamercolor{background canvas}{bg=DPurple!20}
\begin{frame}[plain,noframenumbering]
	
	\begin{center}
		\Huge
		Thanks for listening!
		
		\vspace{1em}
		Any Questions/Comments?
	\end{center}
\end{frame}}
% section questions (end)

%--- Next Frame ---%

\end{document}
